\documentclass[13pt, a4paper]{article}

\usepackage{amsmath}
\usepackage{amssymb}
\usepackage{amsthm}
\usepackage{bussproofs}
\usepackage{fancyhdr}
\usepackage[margin=3cm]{geometry}
\usepackage{hyperref}
\usepackage{mathpartir}
\usepackage{stmaryrd}

\newtheorem{theorem}{Theorem}
\newtheorem{lemma}{Lemma}
\newtheorem{proposition}{Proposition}

\pagestyle{fancy}
\fancyhf{} % Clear all header and footer fields

% Header: left, center, right
\fancyhead[L]{Programming Languages 2}
\fancyhead[C]{Homework 06}
\fancyhead[R]{\today}
\fancyfoot[C]{\thepage}

\begin{document}

Simon Sutoris 7542170

%===============================================================================
\section*{1 STLC with Extensions}
%===============================================================================

\begin{enumerate}

\item % 1
$\texttt{NN} = \texttt{Nat}\times\texttt{Nat}$

\begin{prooftree}
  \AxiomC{$b:\texttt{Bool} \in p : \texttt{NN}, b:\texttt{Bool}$}
  %\RightLabel{\textsc{T-Var}}
  \UnaryInfC{$p : \texttt{NN}, b:\texttt{Bool} \vdash b:\texttt{Bool}$}
  \AxiomC{$p : \texttt{NN} \in p : \texttt{NN}, b:\texttt{Bool}$}
  %\RightLabel{\textsc{T-Var}}
  \UnaryInfC{$p : \texttt{NN}, b:\texttt{Bool} \vdash p : \texttt{NN}$}
  %\RightLabel{\textsc{T-Proj1}}
  \UnaryInfC{$p : \texttt{NN}, b:\texttt{Bool} \vdash p.1 : \texttt{Nat}$}
  %\RightLabel{\textsc{T-IsZero}}
  \UnaryInfC{$p : \texttt{NN}, b:\texttt{Bool} \vdash \texttt{iszero}\,p.1 : \texttt{Bool}$}
  \AxiomC{$p : \texttt{NN} \in p : \texttt{NN}, b:\texttt{Bool}$}
  %\RightLabel{\textsc{T-Var}}
  \UnaryInfC{$p : \texttt{NN}, b:\texttt{Bool} \vdash p : \texttt{NN}$}
  %\RightLabel{\textsc{T-Proj2}}
  \UnaryInfC{$p : \texttt{NN}, b:\texttt{Bool} \vdash p.2 : \texttt{Nat}$}
  %\RightLabel{\textsc{T-IsZero}}
  \UnaryInfC{$p : \texttt{NN}, b:\texttt{Bool} \vdash \texttt{iszero}\,p.2 : \texttt{Bool}$}
  %\RightLabel{\textsc{T-If}}
  \TrinaryInfC{$p : \texttt{NN}, b:\texttt{Bool} \vdash \texttt{if}\,b\,\texttt{then}\,(\texttt{iszero}\,p.1)\,\texttt{else}\,(\texttt{iszero}\,p.2) : \texttt{Bool}$}
  %\RightLabel{\textsc{T-Abs}}
  \UnaryInfC{$p : \texttt{NN} \vdash \lambda b : \texttt{Bool}.\texttt{if}\,b\,\texttt{then}\,(\texttt{iszero}\,p.1)\,\texttt{else}\,(\texttt{iszero}\,p.2): \texttt{Bool} \to \texttt{Bool}$}
\end{prooftree}

\item % 2
$\texttt{NB} = \texttt{Nat}+\texttt{B}$

\begin{prooftree}
  \AxiomC{$x : \texttt{NB} \in x : \texttt{NB}$}
  \UnaryInfC{$x:\texttt{NB}\vdash x : \texttt{NB}$}
  \AxiomC{$y: \texttt{Nat} \in x:\texttt{NB}, y: \texttt{Nat}$}
  \UnaryInfC{$x:\texttt{NB}, y: \texttt{Nat} \vdash y: \texttt{Nat}$}
  \UnaryInfC{$x:\texttt{NB}, y: \texttt{Nat} \vdash \texttt{iszero}\,y : \texttt{Bool}$}
  \AxiomC{$z : \texttt{Bool} \in x:\texttt{NB}, z: \texttt{Bool}$}
  \UnaryInfC{$x:\texttt{NB}, z: \texttt{Bool} \vdash z : \texttt{Bool}$}
  \TrinaryInfC{$x:\texttt{NB}\vdash (\texttt{case}\,x\,\texttt{of}\,\texttt{inl} y \Rightarrow \texttt{iszero}\,y \mid \texttt{inr}\,z \Rightarrow z) : \texttt{Bool}$}
  \UnaryInfC{$\vdash \lambda x : \texttt{NB}.(\texttt{case}\,x\,\texttt{of}\,\texttt{inl} y \Rightarrow \texttt{iszero}\,y \mid \texttt{inr}\,z \Rightarrow z) : \texttt{NB} \to \texttt{Bool}$}
\end{prooftree}

\item % 3
The exercise does not ask for a type derivation, but a term $t : \textbf{A}+(\textbf{A}\times\textbf{A})$:

\begin{prooftree}
  \AxiomC{$\lightning$}
  \UnaryInfC{$x : \textbf{A} \vdash ? : \textbf{A}+(\textbf{A}\times\textbf{A})$}
  \UnaryInfC{$\vdash \lambda x : \textbf{A}.?:\textbf{A}\to (\textbf{A}+(\textbf{A}\times\textbf{A}))$}
\end{prooftree}

Any $t$ with the type $\textbf{A}+(\textbf{A}\times\textbf{A})$ suffice. Two terms I have foud by just plugging stuff in:

\begin{enumerate}
\item $\texttt{inl}\,x : \textbf{A}+(\textbf{A}\times\textbf{A})$
\item $\texttt{inr}\,\{x, x\} : \textbf{A}+(\textbf{A}\times\textbf{A})$
\end{enumerate}

\item % 4
The exercise appears to be malformed. $t_0$ appears to be a record, while the \texttt{case-of} term expects a variant type.

\begin{prooftree}
  \AxiomC{$\lightning$}
  \UnaryInfC{$\vdash \texttt{case}\,\{a = \texttt{true},b = \texttt{false}\}\,\texttt{of} <a=x> \Rightarrow x : ?$}
\end{prooftree}

\item % 5
Derivation:

\begin{prooftree}
  \AxiomC{$p : <c:\texttt{Nat}> \in p : <c:\texttt{Nat}>$}
  \UnaryInfC{$p : <c:\texttt{Nat}> \vdash p : <c:\texttt{Nat}>$}
  \AxiomC{$x: \texttt{Nat} \in p : <c:\texttt{Nat}>,x: \texttt{Nat} $}
  \UnaryInfC{$p : <c:\texttt{Nat}>,x: \texttt{Nat} \vdash x: \texttt{Nat}$}
  \UnaryInfC{$p : <c:\texttt{Nat}>,x: \texttt{Nat} \vdash \texttt{succ}\,x : \texttt{Nat}$}
  \BinaryInfC{$p : <c:\texttt{Nat}> \vdash \texttt{case}\,p\,\texttt{of}\, <c=x> \Rightarrow \texttt{succ}\,x:\texttt{Nat}$}
\end{prooftree}

\end{enumerate}

%===============================================================================
\section*{2 The \texttt{Option} type}
%===============================================================================

\begin{enumerate}

\item % 1

Evaluation rules:

\begin{prooftree}
  \AxiomC{$t \to t'$}
  \RightLabel{\textsc{E-SomeEval}}
  \UnaryInfC{$\texttt{Some}\,t \to \texttt{Some}\,t'$}
\end{prooftree}

\begin{prooftree}
  \AxiomC{$t_1 \to t'_1$}
  \RightLabel{\textsc{E-CaseEval}}
  \UnaryInfC{$\texttt{case}\,t_1\,\texttt{of}\,\texttt{None}\Rightarrow t_2 \mid \texttt{Some}\, x \Rightarrow t_3 \to \texttt{case}\,t'_1\,\texttt{of}\,\texttt{None}\Rightarrow t_2 \mid \texttt{Some}\, x \Rightarrow t_3$}
\end{prooftree}

\begin{prooftree}
  \AxiomC{$$}
  \RightLabel{\textsc{E-CaseNone}}
  \UnaryInfC{$\texttt{case}\,\texttt{None}\,\texttt{of}\,\texttt{None}\Rightarrow t_2 \mid \texttt{Some}\, x \Rightarrow t_3 \to t_2$}
\end{prooftree}

\begin{prooftree}
  \AxiomC{$$}
  \RightLabel{\textsc{E-CaseSome}}
  \UnaryInfC{$\texttt{case}\,\texttt{Some}\,v_1\,\texttt{of}\,\texttt{None}\Rightarrow t_2 \mid \texttt{Some}\, x \Rightarrow t_3 \to [x\mapsto v_1]t_3$}
\end{prooftree}

\item % 2
Typing rules:

\begin{prooftree}
  \AxiomC{$$}
  \RightLabel{\textsc{T-None}}
  \UnaryInfC{$\vdash \texttt{None} : \texttt{Option}\,T$}
\end{prooftree}

\begin{prooftree}
  \AxiomC{$\Gamma \vdash t : T$}
  \RightLabel{\textsc{T-Some}}
  \UnaryInfC{$\Gamma \vdash \texttt{Some}\, t : \texttt{Option}\,T$}
\end{prooftree}

\begin{prooftree}
  \AxiomC{$\Gamma \vdash t_1 : \texttt{Option}\,T_1$}
  \AxiomC{$\Gamma \vdash t_2 : T_2$}
  \AxiomC{$\Gamma, x: T_1 \vdash t_3 : T_2$}
  \RightLabel{\textsc{T-Case}}
  \TrinaryInfC{$\Gamma \vdash \texttt{case}\,t_1\,\texttt{of}\,\texttt{None}\Rightarrow t_2 \mid \texttt{Some}\, x \Rightarrow t_3 : T_2$}
\end{prooftree}

\item % 3

\begin{theorem}\label{thm:2-progress}
Progress: Let $t$ be a closed, well-typed term, that is $\vdash t:T$ for some $T$, then either $t$ is a value or there is some $t'$ with $t\to t'$.
\end{theorem}
\begin{proof}

Straightforward induction on typing derivations. In the cases \textsc{T-None} and \textsc{T-Some}, the term is a value. The only case remaining is \textsc{T-Case}. There exist only 3 cases

\begin{enumerate}
\item $t_1$ is \texttt{None}, in which case \textsc{E-CaseNone} applies
\item $t_1$ is \texttt{Some} $v_1$, in which case \textsc{E-CaseSome} applies and
\item $t_1$ is a reducible term, in which case we can perform inductino over either \textsc{E-SomeEval} or \textsc{E-CaseEval}.
\end{enumerate}

\end{proof}

\begin{theorem}\label{thm:2-progress}
Preservation: If $\Gamma \vdash t:T$ and $t\to t'$, then $\Gamma \vdash t' : T$.
\end{theorem}
\begin{proof}
Induction on a derivation of $t:T$. At each step of the induction, we assume that the desired property holds for all suberivations.

\textit{Case} \textsc{T-None}: $t = \texttt{None}, T = \texttt{Option}\,T$

If the last rule in the derivation is \textsc{T-None}, then we know $t$ must be the constant \texttt{None}, which has the type $\texttt{Option}\,T$. This means $t$ is a value and the theorem doesn't apply.

\textit{Case} \textsc{T-Some}: $t = \texttt{Some}\,x, T = \texttt{Option}\,T$

If the last rule in the derivation is \textsc{T-Some}, then we know $x$ must either be a value or an evaluable term. If $x$ is a value, then $t$ is a value and the theorem does not apply. If $x$ is an evaluable term, \textsc{E-SomeEval} applies. Due to the induction hypothesis, $t$ and $t'$ have the same type, so the requirements for the theorem are satisfied.

\textit{Case} \textsc{T-Case}: $t= \texttt{case}\,t_1\,\texttt{of}\,\texttt{None}\Rightarrow t_2 \mid \texttt{Some}\, x \Rightarrow t_3, T = T_2$

If $t_1$ is an evaluable term, \textsc{E-CaseEval} applies: $t-1 \to t'_1$. Due to te first axiom of \textsc{T-Case}, $t_1$ and $t'_1$ must have the same type and the theorem requirements are satisfied.

If $t_1$ is a value, it's either \texttt{None} or \texttt{Some} $v_1$. In the first, \textsc{E-CaseNone} applies, which due to the second axiom of \textsc{T-Case}, requires $t_2$ to be of type $T_2$, which satisfies the theorem requirements. In the latter, \textsc{E-CaseSome} applies, which due to the third axiom of \textsc{T-Case}, requires $t_3$ to be of type $T_2$, which satisfies the theorem requirements.

\end{proof}

\end{enumerate}

%===============================================================================
\section*{3 Implementation}
%===============================================================================

For my implementation, see \url{https://github.com/Rismosch/Programmiersprachen-2-WiSe-2026/tree/main/typechecker}

%===============================================================================
\section*{4 Substructural type systems}
%===============================================================================

\begin{enumerate}
  
\item % 1
\begin{theorem}\label{thm:4-exchange}
%If $\Gamma_a,x_a:_1T_a,x_b:_1T_b,\Gamma_b \vdash t:T$, then $\Gamma_a,x_b:_1T_b,x_a:_aT_a,\Gamma_b \vdash t:T$.
If $\Gamma_a,x_a:_1T_a,x_b:_1T_b,\Gamma_b \vdash t:T$, then $\Gamma_a,x_b:_1T_b,x_a:_1T_a,\Gamma_b \vdash t:T$.
\end{theorem}
\begin{proof}

We can eliminate common terms using the Context Splitting rules. In the last step we use the comment in the last bullet point, which states:

\begin{quote}
Note that, although we write $\Gamma = \Gamma_1 \circ \Gamma_2$, the order in $\Gamma$ is \textit{not} uniquely determined by $\Gamma_1$ and $\Gamma_2$
\end{quote}

It allows us to reorder terms to produce the same equation on both sides of the equal sign.

\begin{align}
\Gamma_a,x_a:_1T_a,x_b:_1T_b,\Gamma_b &= \Gamma_a,x_b:_1T_b,x_a:_1T_a,\Gamma_b \\
\Gamma_a\circ(x_a:_1T_a,x_b:_1T_b) \circ \Gamma_b &= \Gamma_a\circ(x_b:_1T_b,x_a:_1T_a) \circ \Gamma_b \\
x_a:_1T_a,x_b:_1T_b &= x_b:_1T_b,x_a:_1T_a \\
(x_a:_1T_a) \circ (x_b:_1T_b) &= (x_b:_1T_b) \circ (x_a:_1T_a) \\
(x_a:_1T_a) \circ (x_b:_1T_b) &= (x_a:_1T_a) \circ (x_b:_1T_b)
\end{align}

\end{proof}
  
\item % 2
\begin{theorem}\label{thm:4-exchange}
Weakening does not hold.
\end{theorem}
\begin{proof}
Find $\Gamma, x, t, T, T'$ such that $\Gamma \vdash t : T$ and $x\notin\text{dom}(\Gamma)$ but $\Gamma,x:_1T'\vdash t:T$ does \textit{not} hold.

Chose:

\begin{align}
  \Gamma &= \emptyset\\
  x &= y\\
  t &= z\\
  T &= A\\
  T' &= A\\
\end{align}

The term is well typed:

\begin{prooftree}
  \AxiomC{$$}
  \RightLabel{\textsc{T-Var}}
  \UnaryInfC{$z:_1 A \vdash z : A$}
\end{prooftree}

But using weakening not:

\begin{prooftree}
  \AxiomC{$\lightning$}
  \UnaryInfC{$z:_1 A, y :_1 A \vdash z : A$}
\end{prooftree}

\textsc{T-Var} only applies when the context only includes the variable of the term. If the type system would have the weakening property, then using it would break well-typed terms.

\end{proof}
  
\item % 3
Without weakening, we cannot freely add variables to our context. This implies that the variables in our context represent values that we must use in our term.

\item % 4
This is useful for resource management. A value which is introduced must be consumed somehow.

The closest example from my experience is \texttt{std::thread} in C++. The programmer must call \texttt{std::join} (wait for it to finish its execution) or \texttt{std::detach} (permit the thread to execute independetly), thus consuming it. Linearity is enforced at runtime: If a thread is not consumed, the destructor of \texttt{std::thread} will call \texttt{std::terminate}, which ends the entire program. This means the programmer needs to be explicit about how they use the thread.

Rusts ownership model kinda fits, but not quite. Rust allows implicit dropping, which means variables can go out of scope without explicitly dropping them.

\end{enumerate}

%===============================================================================
\section*{5 Types and Programming Languages}
%===============================================================================

No questions.

\end{document}
